%%%%%%%%%%%%%%%%%%%%%%%%%%%%%%%%%%%%%%%%%%%%%%%%%%%%%%%%%%%%%%%%%%%%%%
% Lớp 12 | Ứng dụng đạo hàm để khảo sát hàm số | Đường tiệm cận
% Dạng 2D1?4-4: Toán thực tế ứng dụng tiệm cận
% Số câu: 12  (sinh tự động bởi tach_cau.py; từ nay thầy sửa trực tiếp file này)
%%%%%%%%%%%%%%%%%%%%%%%%%%%%%%%%%%%%%%%%%%%%%%%%%%%%%%%%%%%%%%%%%%%%%%

%%%%%%%%%%%%%%%% ĐÚNG/SAI %%%%%%%%%%%%%%%%

% nguồn: Phan1_Lop_12_Dai_so_Chuong_1
\begin{ex}%[2D1V4-4]%[ma:5333bab5]
	Cho hàm số $y=\dfrac{x^2-3x+6}{x-1}$ có  đồ thị $(C)$.
	 \choiceTF
		{\True Hàm số có đạo hàm $y'=\dfrac{x^2-2x-3}{\left(x-1\right)^2}$} 
		{Hàm số  có tập xác định $\mathscr{D}=\mathbb{R}\backslash\{-1\}$}
		{Các khoảng nghịch biến của hàm số  chứa $5$ số nguyên}
		{Gọi  $I$ tâm đối xứng của đồ thị $(C)$. Tiếp tuyến của đồ thị tại điểm có hoành độ $x=2$ cắt hai đường tiệm cận tại $A$, $B$. Diện tích tam giác $IAB$ bằng $12$}
	 \loigiai{
		\begin{itemchoice}
			\itemch Ta  có
			\begin{eqnarray*}
				 y'&=&\dfrac{\left(x^2-3x+6\right)'(x-1)-\left(x^2-3x+6\right)(x-1)'}{(x-1)^2} \\
				&=&\dfrac{(2x-3)(x-1)-(x^2-3x+6) \cdot 1}{\left(x-1\right)^2}\\ 
				&=&\dfrac{x^2-2x-3}{\left(x-1\right)^2}. 
			\end{eqnarray*}
			\itemch  Tập xác định của hàm số là $\mathscr{D}=\mathbb{R}\backslash\{1\}$.
			\itemch Ta có $y'=0 \Leftrightarrow x^2-2x-3 = 0 \Leftrightarrow \hoac{&x=-1 \\&x=3.}$ \\ 
			Bảng xét  dấu
			\begin{center}
				\begin{tikzpicture}[scale=.75,transform shape]
					\tkzTabInit[nocadre=false,lgt=1.2,espcl=2.5,deltacl=0.6]{$x$ /1,$f’(x)$  /1}
						{$-\infty$,$-1$,$1$,$3$,$+\infty$} 
					\tkzTabLine{,+,0,-,d,-,0,+,} 
				\end{tikzpicture}
			\end{center}
			Hàm số nghịch biến  trên khoảng $(-1;1)$ và $(1;3)$ chứa các số nguyên là $0$; $2$.
			 \itemch Ta có
			\[\lim\limits_{x \rightarrow 1^+} y= \lim\limits_{x \rightarrow 1^+} \dfrac{x^2-3x+6}{x-1} =  +\infty.\]
			\[\lim\limits_{x \rightarrow 1^-} y=  \lim\limits_{x \rightarrow 1^-} \dfrac{x^2-3x+6}{x-1} = -\infty.\]
			$\Rightarrow x=1$ là đường tiệm cận đứng của đồ thị hàm số. 
			Ta  có
			\[y=\dfrac{x^2-3x+6}{x-1}=\dfrac{(x-1)^2 -(x-1)  +4}{x-1}=x-2+\dfrac{4}{x-1}.\]
			$\Rightarrow y=x-2$ là đường  tiệm cận xiên của đồ thị hàm số. \\
			Tâm đối xứng $I$ của đồ thị $(C)$ là giao điểm của hai đường tiệm cận. \\ 
			 $\Rightarrow I(1;-1)$. \\
			Ta có $y'(2) = -3$; $y(2) =  4$. \\
			Phương trình tiếp tuyến của đồ thị tại điểm  có hoành độ $x=2$ là
			\begin{eqnarray*}
				&&y-4=-3 \cdot (x-2)\\ 
				&\Leftrightarrow&y=-3x+10. 
			\end{eqnarray*}
			\begin{itemize}
				\item Tiếp tuyến  giao với đường tiệm cận đứng $x=1$ tại điểm $A(1;7)$.
				\item Tiếp tuyến giao với đường tiệm  cận xiên tại điểm $B(3;1)$.
			\end{itemize}
			 Ta có $IA=8$, $IB=2\sqrt{2}$, $AB=2\sqrt{10}$ \\
			Nửa chu vi tam giác $IAB$ là  $p=4+\sqrt{2}+\sqrt{10}$. \\
			Diện tích tam giác $IAB$  là
			\begin{eqnarray*}
				S&=&\sqrt{(4+\sqrt{2}+\sqrt{10})(4+\sqrt{2}+\sqrt{10}-8)(4+\sqrt{2}+\sqrt{10}-2\sqrt{2})(4+\sqrt{2}+\sqrt{10}-2\sqrt{10})} \\ 
				&=& 8. 
			\end{eqnarray*}
		\end{itemchoice}
	 }
\end{ex}

% nguồn: Phan1_Lop_12_Dai_so_Chuong_1
\begin{ex}%[2D1V4-4]%[ma:cd49f953]
	Một bể bơi chứa $5\,000$ lít nước tinh khiết. Người ta  bơm vào bể đó nước muối có nồng độ $30$ gam muối cho mỗi lít nước với tốc độ $25$ lít/phút.
	 \choiceTF
		{\True Sau $t$ phút khối lượng muối trong bể là $750t$ (gam)} 
		{Nồng độ muối trong bể sau $t$ phút (tính bằng tỉ số của khối  lượng muối trong bể và thể tích nước trong bể, đơn vị: gam/lít) là $f(t)=\dfrac{30t}{200-t}$}
		{\True Xem $y=f(t)$ là một hàm số xác định trên nửa khoảng $[0;+\infty)$, tiệm cận ngang của đồ thị hàm số  đó có phương trình là $y = 30$}
		{\True Khi $t$ ngày càng lớn thì  nồng độ muối trong bể sẽ tiến gần đến mức $30$ (gam/lít)}
	\loigiai{ 
		\begin{itemchoice}
			\itemch Sau $t$ phút, khối lượng muối trong bể  là
			\[ 
			25\cdot 30\cdot  t=750t \text{ (gam)}
			 \]
			\itemch Thể tích của lượng nước  trong bể là $5\,000+25t$ (lít).\\
			Vậy nồng độ muối sau $t$  phút là
			 \[
			 f(t)=\dfrac{750t}{5\,000+25t}=\dfrac{30t}{200+t}\,\text{ (gam/lít)}.
			 \]
			 \itemch
			 Ta có\\
			$\lim\limits_{t\to +\infty}f(t)=\lim\limits_{t \to +\infty}\dfrac{30t}{200+t}=\lim\limits_{t\to  +\infty}\left(30-\dfrac{6\,000}{200+t}\right)=30$.\\
			Vậy  đường thẳng $y=30$ là tiệm cận ngang của đồ thị hàm số $f(t)$.
			\begin{center}
				\begin{tikzpicture}[scale=.1,xscale=0.1,font=\footnotesize,line join=round,line cap=round,>=stealth]
					\draw[->] (-.5,0)--(0,0) node[below  right]{$O$}--(500,0) node[below]{$x$};
					\draw[->] (0,-1)  --(0,34) node[right]{$y$};
					\draw[blue]  [domain=0:500,samples=100]
						plot  (\x,{(30*(\x))/((\x)+200)});
					\draw[fill]  (0,0) circle (1pt);
					 \foreach \y/\g in {30/180}
					\draw[fill] (0,\y) circle (1pt)  node [shift={(\g:.3)}] {$\y$};
					\draw[red]  (-.1,30)--(500,30);
				\end{tikzpicture}
			\end{center}
			 \itemch
			Ta có đồ thị hàm số $y=f(t)$ nhận đường thẳng $y=30$ làm đường tiệm cận  ngang, tức là khi $t$ càng lớn thì nồng độ muối trong bể sẽ tiến gần đến mức $30$ (gam/lít). Lúc đó, nồng độ muối trong bể sẽ gần như bằng nồng độ nước muối bơm vào bể.
		\end{itemchoice}
	} 
\end{ex}


%%%%%%%%%%%%%%%% TRẢ LỜI NGẮN %%%%%%%%%%%%%%%%

% nguồn: Phan2_Lop_12_Dai_so_Chuong_1
\begin{bt}%[2D1H4-4]%[ma:bfaea089]
	Tại một công ty sản xuất đồ chơi an toàn cho trẻ em, công ty phải chi $34\,567$ USD để thiết lập dây chuyền sản xuất ban đầu. Sau đó, cứ sản xuất được một sản phẩm đồ chơi $X$, công ty phải trả $8$ USD  cho nguyên liệu ban đầu và nhân công. Gọi $x$ ($x \ge 1$) là số đồ chơi $X$ mà công ty đã sản xuất và $C(x)$ (đơn vị: USD) là tổng số tiền bao gồm cả chi phí ban đầu mà công ty phải chi trả khi sản xuất $x$ đồ chơi $X$. Khi đó chi phí trung bình cho mỗi sản phẩm đồ chơi $X$ là hàm số $\overline{C}(x)=\dfrac{C(x)}{x}$ xác định trên $[1;+\infty)$. Khi số đồ chơi sản xuất tăng lên thì chi phí trung bình cho mỗi sản phẩm đồ chơi $X$ giảm xuống nhưng không xuống dưới mức tối thiểu là bao nhiêu USD?
	\shortans{8} 
	 \loigiai{
		Ta có $C(x)  = 34\,567 + 8x$.\\
		Khi đó $\overline{C}(x) =  \dfrac{C(x)}{x} = \dfrac{34\,567 + 8x}{x} = 8 + \dfrac{34\,567}{x}$.\\
		Ta có $\lim\limits_{x \to +\infty} \overline{C}(x) =  \lim\limits_{x \to +\infty} \left(8 + \dfrac{34\,567}{x} \right) = 8$.\\
		Mà $\overline{C}(x)$ là hàm nghịch biến trên $[1;+\infty)$ nên chi phí  trung bình cho mỗi sản phẩm đồ chơi $X$ giảm xuống thấp nhất và không dưới $8$ USD.
	} 
\end{bt}

% nguồn: Phan2_Lop_12_Dai_so_Chuong_1
\begin{bt}%[2D1H4-4]%[ma:8fcb64d2]
	Số lượng sản phẩm bán được của một công ty trong $x$ (tháng) được  tính theo công thức $S(x)=200\left(5-\dfrac{9}{2+x}\right)$, trong đó $x \ge 1$. Số lượng sản phẩm lớn nhất mà công ty có thể bán được trong $x$ tháng là
	\shortans{1000} 
	 \loigiai{
		Số lượng sản phẩm bán được lớn nhất của công ty  là tiệm cận ngang của hàm số $S(x)$. \\
		Ta có $\lim\limits_{x  \to +\infty} y = \lim\limits_{x \to +\infty} 200\left(5-\dfrac{9}{2+x}\right)=1\,000$.\\
		 Vậy đường tiệm cận ngang của đồ thị hàm số $S(x)$ là $y=1\,000$.
	} 
\end{bt}

% nguồn: Phan2_Lop_12_Dai_so_Chuong_1
\begin{bt}%[2D1H4-4]%[ma:fdb9ead0]
	Một bể chứa $5\,000$ lít nước tinh khiết. Người ta bơm vào bể đó nước muối có nồng độ $30$ gam muối cho mỗi lít nước với tốc độ $25$ lít/phút. Nồng độ muối trong bể được tính bởi công thức $f(t) = \dfrac{30t}{200+t}$ ($t$ là thời gian nước chảy vào bể (phút)). Nồng độ muối lớn nhất mà bể có  thể đạt được là bao nhiêu gam/lít?
	 \shortans{30}
	\loigiai{ 
		Nồng độ muối lớn nhất mà bể có thể  đạt được là tiệm cận ngang của đồ thị hàm số. \\
		Ta  có
		$\lim\limits_{t  \to +\infty} f(t) = \lim\limits_{t \to +\infty} \dfrac{30t}{200+t} = \lim\limits_{t \to +\infty} \left(30 - \dfrac{6000}{200+t}\right) = 30$. \\
		Vậy nồng độ  muối lớn nhất của bể là $30$ gam/lit.
	} 
\end{bt}

% nguồn: Phan2_Lop_12_Dai_so_Chuong_1
\begin{bt}%[2D1V4-4]%[ma:1695c4d5]
	Một công ty sản xuất máy tính đã xác định được rằng, tính trung bình một nhân viên có thể lắp ráp được $N(x)=\dfrac{50 x}{x+4}(x \geq 0)$ bộ phận mỗi ngày sau $x$ ngày đào tạo. Xem $y=N(x)$ là một hàm  số xác định trên $[0;+\infty)$, khi số ngày đào tạo tăng lên, hãy tính số bộ phận một nhân viên lắp ráp tối đa không vượt quá bao nhiêu?
	 \shortans{50}
	 \loigiai{
		Ta có $\displaystyle\lim  _{x \rightarrow+\infty} N(x)=\displaystyle\lim _{x \rightarrow+\infty} \frac{50 x}{x+4}=\displaystyle\lim _{x \rightarrow+\infty} \frac{50}{1+\frac{4}{x}}=50$.
		Vậy một nhân viên lắp ráp tối đa không  vượt quá $50$ bộ phận máy tính.
	} 
\end{bt}

% nguồn: Phan2_Lop_12_Dai_so_Chuong_1
\begin{bt}%[2D1V4-4]%[ma:1be55739]
	Nồng độ thuốc trong máu của một bệnh nhân sau khi tiêm $t$ giây là $P(t) \left(\dfrac{mg}{m l}\right)$. Với $P(t)$ cho bởi công thức $P(t)=\sqrt{0{,}09 t^{2}+0{,}018t+0{,}05}-0{,}3t$. Nồng độ thuốc tồn dư, tức nồng độ thuốc vẫn còn trong cơ thể bệnh  nhân trong dài hạn khi tiêm là bao nhiêu?
	 \shortans{0,03}
	\loigiai{ 
		Do đề bài cho  nồng độ thuốc trong cơ thể bệnh nhân trong dài hạn nên $t \to +\infty$.\\
		Ta  có
		\begin{eqnarray*}
			\lim\limits_{t \to +\infty} P(t)&=&\lim\limits_{t \to +\infty}\left(\sqrt{0{,}09 t^{2}+0{,}018 t+0{,}05}-0{,}3 t\right)  \\
			&=&\lim\limits_{t  \to+\infty} \dfrac{0{,}09 t^{2}+0{,}018 t+0{,}05-(0{,}3 t)^{2}}{\sqrt{0{,}09 t^{2}+0{,}018 t+0{,}05}+0{,}3 t} \\
			&=&\lim\limits_{t \to+\infty} \dfrac{0{,}018 t+0{,}05}{\sqrt{0{,}09 t^{2}+0{,}018 t+0{,}05}+0{,}3 t}  \\
			&=&\lim\limits_{t \to +\infty}  \dfrac{t\left(0{,}018+\dfrac{0{,}05}{t}\right)}{t\left(\sqrt{0{,}09+\dfrac{0{,}018}{t}+\dfrac{0{,}05}{t^{2}}}+0{,}3\right)} \\
			&=&\lim\limits_{t  \to +\infty} \dfrac{0{,}018+\dfrac{0{,}05}{t}}{\sqrt{0{,}09+\dfrac{0{,}018}{t}+\dfrac{0{,}05}{t^{2}}}+0{,}3}=\dfrac{0{,}018}{\sqrt{0{,}09}+0{,}3}=0{,}03.
		\end{eqnarray*}
		Vậy nồng độ thuốc tồn dư trong cơ thể bệnh nhân  là $0{,}03\left(\dfrac{\mathrm{mg}}{\mathrm{ml}}\right)$.
	} 
\end{bt}

% nguồn: Phan2_Lop_12_Dai_so_Chuong_1
\begin{bt}%[2D1V4-4]%[ma:c103a57d]
	Giả sử doanh số $f(t)$ (tính bằng số sản phẩm) của một sản phẩm mới (trong vòng một số năm nhất định) tuân theo quy luật logistic được mô  hình hoá bằng hàm số $f(t) = \dfrac{5000}{1 + 5e^{-t}}, (t \ge 0)$, trong đó thời gian $t$ được tính bằng năm kể từ khi phát hành sản phẩm mới. Khi đó đạo hàm $f'(t)$ là biểu thị tốc độ bán hàng. Hỏi sau khi phát hành bao~nhiêu năm thì tốc độ bán hàng là lớn nhất? (làm tròn kết quả đến hàng phần chục).
	\shortans{1,6} 
	\loigiai{ 
		Xét hàm  số:
		$$ f(t) = \dfrac{5000}{1+5e^{-t}}, (t \ge 0)  $$
		Vì đạo hàm $f'(t)$ là biểu thị tốc độ bán hàng nên để biết sau khi phát hành bao~nhiêu năm thì tốc độ bán hàng là lớn nhất ta  tìm GTLN của hàm số:
		$$ f'(t) = \dfrac{25000e^{-t}}{(1+5e^{-t})^2}, (t \ge 0)  $$
		$$  f''(t) = \dfrac{25000e^{-t}(5e^{-t}-1)}{(1+5e^{-t})^3} $$
		$$ f''(t)=0 \Leftrightarrow \dfrac{25000e^{-t}(5e^{-t}-1)}{(1+5e^{-t})^3} = 0 \Leftrightarrow  5e^{-t}-1=0 \Leftrightarrow e^{-t}=\dfrac{1}{5} \Leftrightarrow t=\ln 5 $$
		Bảng  biến thiên:
		\begin{center}
			\begin{tikzpicture}[scale=0.7,transform shape]
				 \tkzTabInit[nocadre,espcl=5]
					{$t$/1,$f'(t)$/1,$f  (t)$/3}
					{$0$,$ \ln 5 $,$ +\infty  $}
				\tkzTabLine{,+,0,-,} 
				 \tkzTabVar{-/,+/$\max$,-/}
			\end{tikzpicture}
		\end{center}
		 Hàm số $f'(t)$ đạt GTLN khi $t=\ln 5 \approx 1,6$. \\
		Vậy  sau khi phát hành 1,6 năm thì tốc độ bán hàng là lớn nhất.
	 }
\end{bt}

% nguồn: Phan2_Lop_12_Dai_so_Chuong_1
\begin{bt}%[2D1V4-4]%[ma:8bdcac8c]
	Chi phí xuất bản $x$ cuốn tạp chí (bao gồm: lương cán bộ, công nhân viên, giấy in...) được cho bởi $C(x) = x^2 - 2000x + 10^8$ đồng. Chi phí phát hành cho mỗi cuốn là 4 nghìn đồng ($4000$ đồng). $M(x) = \dfrac{T(x)}{x}$ với $T(x)$ là tổng chi phí (xuất bản và phát hành) cho $x$  cuốn tạp chí, được gọi là chi phí trung bình cho một cuốn tạp chí khi xuất bản $x$ cuốn. Khi số lượng cuốn tạp chí phát hành cực lớn thì chi phí trung bình cho mỗi cuốn tạp chí $M(x)$ sẽ tiệm cận với đường thẳng có phương trình dạng $y=ax+b$. Tính $P=68a+3b+800$.
	 \shortans{6868}
	 \loigiai{
		Theo giả thiết, ta có: \\ 
		$T(x)=C(x)+4000x = x^2 - 2000x + 10^8 + 4000x = x^2 + 2000x + 10^8$ (đồng). \\ 
		$M(x)  = \dfrac{T(x)}{x} = \dfrac{x^2 + 2000x + 10^8}{x} = x + \dfrac{10^8}{x}+2000 $. \\
		$\lim\limits_{x \to +\infty} [M(x) - (x+2000)] =  \lim\limits_{x \to +\infty} \left[x + 2000 + \dfrac{10^8}{x} - (x+2000) \right] = \lim\limits_{x \to +\infty} \dfrac{10^8}{x} = 0$. \\
		Khi đó đồ thị hàm số $M(x) =  x + 2000 + \dfrac{10^8}{x}$ có 1 tiệm cận xiên là $y = x+2000$. \\
		Khi số lượng cuốn tạp chí phát hành  cực lớn thì chi phí trung bình cho mỗi cuốn tạp chí $M(x)$ sẽ tiệm cận với đường thẳng $y=x+2000$. \\
		Suy ra $\heva{&a=1  \\ &b=2000.}$ \\
		Vậy $P = 68a +  3b + 800 = 68 \cdot 1 + 3 \cdot 2000 + 800 = 6868$.
	 }
\end{bt}

% nguồn: Phan2_Lop_12_Dai_so_Chuong_1
\begin{bt}%[2D1V4-4]%[ma:996158d8]
	Một bể chứa $1\,000$ lít nước tinh khiết. Người ta bơm vào bể đó nước muối có nồng độ $15$ gam muối cho mỗi lít nước  với tốc độ $20$ lít/phút. Biết rằng nồng độ muối trong bể sau $t$ phút (tính bằng tỉ số của khối lượng muối trong bể và thể tích nước trong bể, đơn vị: gam/lít) là một hàm số $f(t)$, thời gian $t$ tính bằng phút. Phương trình tiệm cận ngang của đồ thị hàm số $f(t)$ bằng là $y=a$, $(a \in \mathbb{R})$. Tìm $a$.
	 \shortans{15}
	 \loigiai{
		Sau $t$ phút,  ta có khối lượng muối trong bể là $20 \cdot 15t = 300t$ (gam). \\
		Thế tích của lượng nước trong  bể sau $t$ phút là $1\,000+20t$ (lít). \\
		Do đó nồng độ muối sau  $t$ phút là $f(t) = \dfrac{300t}{1\,000+20t} = \dfrac{30t}{100+2t}$ (gam/lít). \\
		Ta có $\lim\limits_{t \to +\infty} f(t) = \lim\limits_{t \to +\infty} \dfrac{30t}{100+2t} = \dfrac{30}{2} = 15$ nên đồ thị  hàm số $y = f(t)$ có phương trình tiệm cận ngang là $y=15$. \\
		 Vậy $a=15$.
	 }
\end{bt}

% nguồn: Phan2_Lop_12_Dai_so_Chuong_1
\begin{bt}%[2D1V4-4]%[ma:bfc97649]
	Một cơ sở sản xuất tính toán rằng số sản phẩm trung bình mà một nhân viên làm được mỗi ngày là $f(x)=\dfrac{100 x}{x+10}$ với $x$ là số ngày kinh nghiệm làm việc $(x \geq 0)$. Xem $y=f(x)$ là một hàm số xác định trên $[0;+\infty)$. Khi số ngày kinh nghiệm làm việc tăng lên thì số sản  phẩm trung bình tối đa mà một nhân viên có thể làm được trong một ngày là bao nhiêu?
	 \shortans{99}
	 \loigiai{
		Ta có: $f(x)=\dfrac{100 x}{x+10}=100-\dfrac{1000}{x+10}<100$ và $\displaystyle\lim _{x \rightarrow+\infty} f(x)=100$.\\ Nên số sản phẩm tối đa mà một nhân viên  làm được trong một ngày là không quá 99 sản phẩm.
		Ta có $f(x)=99  \Leftrightarrow \dfrac{100 x}{x+10}=99 \Leftrightarrow x=990$.
		Vậy khi số ngày kinh  nghiệm làm việc tăng lên thì số sản phẩm trung bình tối đa mà một nhân viên có thể làm được trong một ngày là $99$ .
	} 
\end{bt}

% nguồn: Phan2_Lop_12_Dai_so_Chuong_1
\begin{bt}%[2D1V4-4]%[ma:c91cf04a]
	Để loại bỏ $x$\,\% chất gây  ô nhiễm không khí từ khí thải của một nhà máy, người ta ước tính chi phí cần bỏ ra là $C(x) = \dfrac{300x}{100 - x}$ (triệu đồng), $0 \le x < 100$. Tìm số đường tiệm cận đứng của đồ thị hàm số $y = C(x)$.
	 \shortans{1}
	\loigiai{ 
		Tập xác định  $\mathscr{D} = [0;100)$. \\
		Xét hàm số $y = C(x) = \dfrac{300x}{100-x}$, $0 \leq x <  100$. \\
		Ta có  $y' = \dfrac{30\,000}{(100-x)^2} > 0$, với mọi $x \in [0;100)$. \\
		Do đó hàm số luôn đồng biến trên  nửa khoảng $[0;100)$. \\
		Vì $\lim\limits_{x \to 100^-} C(x) = \lim\limits_{x \to 100^-} \dfrac{300x}{100-x} = +\infty$, nên đồ thị hàm số có tiệm cận  đứng là $x = 100$.
	 }
\end{bt}
